\begin{remark}[the cusp is not the torsion locus]
At $p_0$, if
$Q=\nu_*\mathcal O_{\mathrm{dP}_6}/\mathcal O_W$, then
\[
 0\longrightarrow\Omega^1_{N_0/\Delta}\longrightarrow
 \mathcal E|_{N_0}\longrightarrow i_*Q\longrightarrow0,
 \qquad h^\bullet(Q)=(2,1,0).
\tag{11.17a}
\]
This relative sequence is the right column of the following exact
normalization--residue diagram; all sheaves in the bottom row are regarded as
sheaves on $N_0$ by $i_*$:
\[
\begin{array}{ccccccccc}
&&0&&0&&0&&\\
&&\big\downarrow&&\big\downarrow&&\big\downarrow&&\\
0&\longrightarrow&f^*\Omega^1_\Delta&\longrightarrow&\Omega^1_{N_0}
 &\longrightarrow&\Omega^1_{N_0/\Delta}&\longrightarrow&0\\
&&\big\downarrow&&\big\downarrow&&\big\downarrow&&\\
0&\longrightarrow&f^*\Omega^1_\Delta(\log0)&\longrightarrow&
 \Omega^1_{N_0}(\log W)&\longrightarrow&\mathcal E|_{N_0}
 &\longrightarrow&0\\
&&\big\downarrow\,{\scriptstyle\operatorname{res}_B}
&&\big\downarrow\,{\scriptstyle\operatorname{res}_W}
&&\big\downarrow\,{\scriptstyle\rho}&&\\
0&\longrightarrow&i_*\mathcal O_W&\longrightarrow&
 i_*\nu_*\mathcal O_{\mathrm{dP}_6}&\longrightarrow&i_*Q
 &\longrightarrow&0\\
&&\big\downarrow&&\big\downarrow&&\big\downarrow&&\\
&&0&&0&&0&&
\end{array}
\tag{11.17a$'$}
\]
The first two columns are the base and normalized-divisor residue sequences;
the bottom row is the definition of $Q$.  Hence the snake lemma identifies
the right column with (11.17a), including its displayed map $\rho$.

The exactness assertion is also visible stalkwise.  At a double chart $t=xy$, with one
coordinate along the double stratum and logarithmic generator $e$, the image
of ordinary relative differentials in $\mathcal E$ contains the stratum
differential and has transverse part $(x,y)e$; its cokernel is therefore
$R/(x,y)$.  At a triple chart $t=xyz$, write
$\mathcal E=Re_1\oplus Re_2$ with
$d\log x=e_1$, $d\log y=e_2$, and
$d\log z=-e_1-e_2$.  The image is generated by
\[
 xe_1,\qquad ye_2,\qquad z(e_1+e_2),
\]
Writing a normalized branch value as a triple modulo the diagonal, the
quotient map is explicitly
\[
 e_1\longmapsto[(1,0,0)],\qquad
 e_2\longmapsto[(0,1,0)].
\tag{11.17a$''$}
\]
Its kernel is exactly
$\langle xe_1,ye_2,z(e_1+e_2)\rangle$, so this is the standard local
presentation of the normalization quotient $Q$.
Away from the double locus the map is an isomorphism.  This proves (11.17a)
stalkwise.
Here are the maps, rather than only their ranks.  Quotienting the normalization
complex (9.10) by $\mathcal O_W$ gives
\[
 0\longrightarrow Q\longrightarrow
 \bigoplus_{k=1}^3\mathcal O_{D_k}
 \xrightarrow{\delta_1}\mathbb C_P\oplus\mathbb C_Q
 \longrightarrow0.
\]
Because $D_k\simeq\mathbb P^1$, its cohomology is the cohomology of
\[
 \mathbb C^3\longrightarrow\mathbb C^2,\qquad
 (a_1,a_2,a_3)\longmapsto(s,s),\qquad s=a_1-a_2+a_3.
\]
Thus
\[
 H^0(Q)=\{(a_1,a_2,a_3):a_1-a_2+a_3=0\}\simeq\mathbb C^2,
 \qquad H^1(Q)=\mathbb C^2/\mathbb C(1,1)\simeq\mathbb C,
\tag{11.17b}
\]
and $H^2(Q)=0$.  Since $\mathrm{dP}_6$ is rational, the connecting maps for
$0\to\mathcal O_W\to\nu_*\mathcal O_{\mathrm{dP}_6}\to Q\to0$ are
isomorphisms
\[
 \delta_0:H^0(Q)\xrightarrow{\sim}H^1(W,\mathcal O_W),
 \qquad
 \delta_1:H^1(Q)\xrightarrow{\sim}H^2(W,\mathcal O_W).
\tag{11.17c}
\]

At a triple chart $t=z_0z_1z_2$, a relative logarithmic form is represented
by $\eta=\sum a_i\,d\log z_i$, with holomorphic coefficient triple
$(a_0,a_1,a_2)$ taken modulo $\mathcal O(1,1,1)$.  Restricting the branch
differences to the three double strata gives
\[
 \rho(\eta)=(a_0-a_1,\ a_0-a_2,\ a_1-a_2).
\tag{11.17d}
\]
Formula (11.17d) is written in the $P$-ordering of the globally labelled
double curves.  The $Q$-point sign is fixed by a separate calculation, not by
an implicit relabelling.  In the common $e_1,e_2$-plane, use the ordered
branch vertices
\[
 (b_0,b_1,b_2)=(0,e_1,e_2),\qquad
 (c_0,c_1,c_2)=(e_1,e_1+e_2,e_2)
\]
at $P$ and $Q$, respectively.  The coefficient triples of
$d\log x_1,d\log x_2$ are
\[
 (0,1,0),\ (0,0,1)\quad\text{at }P,
 \qquad
 (1,1,0),\ (0,1,1)\quad\text{at }Q.
\]
The globally labelled branch-difference maps are
\[
 d_P(a)=(a_0-a_1,a_0-a_2,a_1-a_2),
 \qquad
 d_Q(a)=(a_2-a_1,a_0-a_1,a_0-a_2).
\tag{11.17d$'$}
\]
Direct substitution gives, at both triple points, the same two global
sections
\[
 q_1=(-1,0,1),\qquad q_2=(0,-1,-1).
\tag{11.17d$''$}
\]
For constant coefficients this induces an isomorphism
$\mathbb C^3/\mathbb C(1,1,1)\to H^0(Q)$.  In the toric chart
$(z_0,z_1,z_2)=(t/(x_1x_2),x_1,x_2)$ it sends
\[
 d\log x_1\longmapsto(-1,0,1),\qquad
 d\log x_2\longmapsto(0,-1,-1),
\]
which is a basis of the hyperplane in (11.17b); unimodular changes of cusp
chart merely change this basis by $\mathrm{GL}_2(\mathbb Z)$.

Write $\eta_1=d\log x_1$, $\eta_2=d\log x_2$, take the $q_i$ of
(11.17d$''$), and put $\alpha_i=\delta_0(q_i)$.  By (11.17c),
$\alpha_1,\alpha_2$ form a basis of $H^1(W,\mathcal O_W)$.  Naturality of
the connecting maps fixes a sign that must be retained.  With the standard
\v Cech differential
$(\delta c)_{i_0\ldots i_{r+1}}
=\sum_k(-1)^kc_{i_0\ldots\widehat{i_k}\ldots i_{r+1}}$, lift $q_i$ to the
normalization complex.  The total differential of a degree-one cocycle
$\alpha$ times that lift is
$(-1)^1\alpha\smile\delta\widetilde q_i$.  The projection formula therefore
gives exactly
\[
 \delta_1\bigl(\rho_*(\alpha\otimes\eta_i)\bigr)
 =-\alpha\smile\delta_0(q_i)=-\alpha\smile\alpha_i.
\]
The cup product in \cref{lem:cusp-structure-sheaf-audit} is an isomorphism.
Define, with this convention,
\[
 \theta=-\delta_1^{-1}(\alpha_1\smile\alpha_2).
\]
Graded commutativity now gives, with no residual sign ambiguity,
\[
 \rho_*=(0,1,-1,0):
 \mathbb C\langle
 \alpha_1\eta_1,\alpha_1\eta_2,\alpha_2\eta_1,\alpha_2\eta_2
 \rangle\longrightarrow\mathbb C\theta.
\tag{11.17e}
\]
In particular the map is surjective.  Applying direct images to (11.17a),
the first map is reduction modulo $t$ followed by the isomorphism (11.17d),
and the second is (11.17e).  Consequently
\[
 f_*\Omega^1_{N_0/\Delta}=t\mathcal O^2,\quad
 R^1f_*\Omega^1_{N_0/\Delta}=\ker(\mathcal O^4\to\mathbb C_{p_0}),
 \quad R^2f_*\Omega^1_{N_0/\Delta}=\mathcal O^2.
\]
For the row in (11.17e), the middle kernel is generated by
$e_1,e_4,e_2+e_3,te_2$.
Thus $\Omega^1_{X/B}$ is torsion-free in a neighbourhood of the cusp, but
non-locally-free along its $xy$ and $xyz$ strata.  Globally its torsion is
supported precisely on the two multiple fibres, as the following local module
calculation proves.  On the smooth cover $J'_j$, let
$R=\mathcal O_{J'_j,x}$ and choose coordinates $(s,y_1,y_2)$ smooth over the
$s$-disc.  Relative to the original base $B$, where $t=s^{m_j}$, one has
\[
 \Omega^1_{J'_j/B,x}
 =\frac{R\,ds\oplus R\,dy_1\oplus R\,dy_2}
        {(m_js^{m_j-1}ds)}
 \simeq R\,dy_1\oplus R\,dy_2
 \oplus\bigl(R/(s^{m_j-1})\bigr)ds.
\tag{11.17f}
\]
Thus the last summand is exactly the torsion submodule.  The quotient map
$q:J'_j\to N_j$ is finite \'{e}tale, so the relative cotangent sequence gives
\[
 q^*\Omega^1_{N_j/B}\simeq\Omega^1_{J'_j/B}.
\]
Faithful flatness descends both the presence of the displayed torsion and its
absence off $s=0$.  Hence ordinary relative torsion is supported exactly on
the two multiple fibres and is absent at the cusp.  The reconstructed
manuscript's assertion that torsion occurs at all three singular fibres is
false.
\end{remark}
